\section{Introduction}
\label{sec:intro}

Images captured in low light suffer from poor visibility, strong sensor noise, colour distortion
and loss of detail. A camera's image signal processor (ISP) turns the raw sensor data into an
sRGB image through denoising, demosaicing, white balance, colour correction and tone mapping.
When the photon count is very low, these steps fail: brightening the image also amplifies the
noise, and denoisers remove texture together with noise.

Learned ISPs map short-exposure raw images directly to long-exposure references. Early methods
are regression networks trained with pixel losses~\cite{chen2018sid,wei2020eld}. They tend to
oversmooth regions where the signal is very weak. Diffusion models~\cite{ho2020ddpm} can generate
plausible detail instead of averaging it away. DarkDiff~\cite{zheng2025darkdiff} goes further: it
retasks a \emph{pretrained} Stable Diffusion model~\cite{rombach2022ldm} as the ISP, so the
strong natural-image prior of text-to-image pretraining is reused for restoration. It reports
the best perceptual quality on the SID~\cite{chen2018sid}, ELD~\cite{wei2020eld} and
LRD~\cite{zhang2023lrd} benchmarks.

This quality has a cost. DarkDiff fine-tunes the full \sdtwo{} U-Net (866M parameters), plus a
copy of its encoder used as a context processor, for 3000 epochs on eight A100 GPUs. Inference
runs 50 denoising steps, each with two U-Net passes for classifier-free guidance. This puts the
method out of reach of modest hardware.

Our goal is a low-light enhancement model that keeps the benefits of a pretrained diffusion
prior but runs on smaller GPUs. We make three contributions:
\begin{enumerate}
  \item We swap the \sdtwo{} U-Net for three distilled U-Nets that are pretrained in the same
        latent space: \bksdm{}~\cite{kim2024bksdm}, \tinysd{}~\cite{segmind2023tinysd} and
        \sdxs{}~\cite{song2024sdxs}. Each has about 330M parameters, 62\% fewer than \sdtwo{}.
  \item We add \emph{illumination-aware conditioning}. A low-resolution illumination map of the
        input modulates the condition features, biases the attention keys and gates the output of
        every region cross-attention layer. It lets the model treat dark, low-SNR regions
        differently from bright ones. All new layers start at zero, so the model starts as plain
        DarkDiff.
  \item We re-implement DarkDiff in PyTorch and compare the original model with our compact,
        illumination-guided variants at an equal training budget. We report image quality
        (PSNR, SSIM, LPIPS) and efficiency (parameters, FLOPs, time and memory).
\end{enumerate}
